\chapter{The System Browser}
\label{ch:browser}

\standfirst{One window, five panes and six menus, and every line of the
system behind them. This is where the work happens, and the half hour spent
learning its menus is the best-paid half hour in Smalltalk.}

\section{The shape of it}

\ct{browser} on the desktop menu sends \ct{BrowserView openOn: SystemOrganization}
--- a browser on the whole system. Four list panes across the top, a code pane
across the bottom, and two switches under the second list.

\begin{center}
\small
\begin{tabular}{@{}lp{0.58\textwidth}@{}}
\toprule
1. category & groups of classes, e.g.\ \ctp{Collections-Unordered}\\
2. class & the classes in that group\\
3. message category & groups of methods, e.g.\ \ctp{accessing}\\
4. message & the selectors in that group\\
\midrule
code pane & whatever the four above have narrowed down to\\
\bottomrule
\end{tabular}
\end{center}

Selection runs left to right: choosing a category fills the class list,
choosing a class fills the message categories, choosing one of those fills the
messages. Deselecting a pane empties everything to its right. There is no
other state --- what you are looking at is exactly what the four selections
say you are looking at.

The wheel scrolls any of the five panes, this system having taught
\ct{ScrollController} and \ct{ParagraphEditor} what a wheel is; 1983 had no
such thing and scrolled with a bar on the left edge that appears when the
pointer enters the pane.

\begin{stnote}[Source comes from a file]
The image does not store method source. It stores a pointer into the sources
or changes file and reads the text back when you select a method, which is why
Chapter~\ref{ch:keeping} is so insistent that the changes file travels with the
image. A browser showing you nothing for a method you wrote is that file
missing, not the method.
\end{stnote}

\section{\ctp{instance} and \ctp{class}}

Under the class list are two switches. \ct{instance} shows the methods an
instance of the selected class understands; \ct{class} shows the methods the
class object itself understands --- \ct{new}, and the rest of
Chapter~\ref{ch:classside}. They are the same browser looking at two different
objects, and forgetting which one is lit is the most common way to write a
method that is never called.

\section{What the code pane is showing}

The code pane has no scrollback and no tabs. It shows one thing, and which
thing depends on what you last chose:

\begin{center}
\small
\begin{tabular}{@{}lp{0.58\textwidth}@{}}
\toprule
a message selected & that method's source\\
a class, no message category & its \emph{definition} --- the
  \ctp{subclass:} \ldots{} \ctp{category:} expression\\
\ctp{comment} chosen & the class comment\\
\ctp{hierarchy} chosen & the class's ancestors and descendants, as text\\
\ctp{protocols} chosen & the class's message categories, as editable text\\
\ctp{edit all} chosen & the whole category list, as editable text\\
\bottomrule
\end{tabular}
\end{center}

The last four are modes you enter from a menu and leave by selecting something
else. They are not separate windows, which is worth knowing the first time
\ct{comment} appears to have replaced your method.

\section{Accept is the compiler}

Editing in the code pane changes nothing. \ct{accept} on its middle-button
menu compiles what is there and installs it, and the installation is immediate
--- no build, nothing to restart, no possibility of a stale version running.

What \ct{accept} means depends on the mode above. On a method it compiles a
method into the selected class, on the selected side. On a class definition it
redefines the class and recompiles every method affected by the change. On a
comment it stores the comment. On \ct{protocols} or \ct{edit all} it
reorganizes.

\begin{stwarn}
A code pane you have edited and not accepted holds the only copy of that text.
\ct{cancel} throws it away without asking, and choosing another item in any
list pane asks first --- that question is the system's one guard against
losing work, so read it rather than dismissing it.
\end{stwarn}

\section{The category pane menu}

With a category selected:

\begin{description}
\item[\ctp{file out}] Write every class in the category to
  \ctp{Category-Name.st}, in superclass order, in chunk format.
\item[\ctp{print out}] The same, formatted for reading rather than for filing
  back in.
\item[\ctp{spawn}] A new browser on this category alone.
\item[\ctp{add category}] Prompt for a name and insert it.
\item[\ctp{rename}] Prompt for a new name and rename it in place.
\item[\ctp{remove}] Remove the category \emph{and every class in it}. It asks
  ``are you certain that you want to remove all classes in this category?''
  first, and that question is not a formality.
\item[\ctp{update}] Re-read the category list from \ct{SystemOrganization}.
  What you want after filing code in from somewhere else, or after another
  process has added a class.
\item[\ctp{edit all}] The whole category list as text in the code pane.
  Reorder it, add to it, \ct{accept}.
\end{description}

With no category selected the menu shrinks to \ct{add category},
\ct{update} and \ct{edit all}, those being the three that do not need one.

\section{The class pane menu}

Thirteen items in four groups, and the middle two groups are the ones worth
learning.

\begin{description}
\item[\ctp{file out}, \ctp{print out}] As above, for this class: both sides,
  one file.
\item[\ctp{spawn}] A browser on this class alone.
\item[\ctp{spawn hierarchy}] A browser whose category list holds one entry,
  \ct{**Hierarchy**}, and whose class list is this class's ancestors followed
  by all of its descendants. The best way there is to read an inherited method
  and its overrides together.
\medskip
\item[\ctp{hierarchy}] The same information as text in the code pane, without
  opening a window.
\item[\ctp{definition}] The \ctp{subclass:} expression. Edit it and
  \ct{accept} to add an instance variable, change the superclass or move the
  class to another category; every affected method is recompiled for you.
\item[\ctp{comment}] The class comment, editable. Write one.
\item[\ctp{protocols}] The class's message categories as editable text.
\medskip
\item[\ctp{inst var refs}] Pops up a menu of the instance variables of this
  class \emph{and its superclasses}; choose one and get a browser of every
  method that reads or writes it. This is how you find out what a variable is
  for.
\item[\ctp{class var refs}] The same for class variables.
\item[\ctp{class refs}] Every method anywhere in the system that mentions this
  class by name. How you find out who uses what you are about to change.
\medskip
\item[\ctp{rename}] Rename the class. References to it are not rewritten ---
  \ctp{class refs} first.
\item[\ctp{remove}] Remove it from the system.
\end{description}

\section{The message category pane menu}

The third pane has a menu of its own, which is easy to go years without
noticing:

\begin{description}
\item[\ctp{file out}, \ctp{print out}] Just the methods in this protocol.
\item[\ctp{spawn}] A browser on this protocol alone.
\item[\ctp{add protocol}] Prompt for a name and add it.
\item[\ctp{rename}, \ctp{remove}] As they read. Removing a protocol removes
  the methods in it.
\end{description}

With nothing selected in that pane the menu is the single item
\ct{add protocol}. Categorising methods is not decoration: it is the only
index a reader of your class has, and \ctp{as yet unclassified} is a
confession.

\section{The message pane menu, and the three that matter}

\begin{description}
\item[\ctp{file out}, \ctp{print out}] This one method.
\item[\ctp{spawn}] A browser on this method alone --- useful when you want to
  keep it on screen while you navigate elsewhere.
\item[\ctp{move}] Prompt for a destination protocol. Type a protocol name to
  move the method within the class; type \ct{Class>protocol} and it is
  \emph{copied} to that protocol of that class instead. The prompt says so, in
  four words, and it is the only place the system explains it.
\item[\ctp{remove}] Remove the method.
\end{description}

The other three are the ones that make a system of five thousand methods
navigable, and they are worth learning by name.

\begin{description}
\item[\ctp{senders}] Every method in the system that sends this selector.
  \ct{Smalltalk browseAllCallsOn:} under the hood. This is how you find out
  who calls something, and therefore whether it is safe to change.
\item[\ctp{implementors}] Every class that implements this selector.
  This is how you find out what a message \emph{means} when the receiver could
  be anything.
\item[\ctp{messages}] A menu of every selector this method sends; choose one
  and you get its implementors. This is how you walk forwards through code you
  are reading, one send at a time.
\end{description}

\begin{stnote}
\ctp{implementors} is the answer to ``what does \ctp{do:} do?'' --- there are
dozens, one per collection, and reading three of them tells you more than any
documentation could. \ctp{senders} is the answer to ``is anyone relying on
this?'' Between them they replace grep, and they are better than grep because
they know the difference between a selector and a word.
\end{stnote}

All three open a \emph{message-list browser}: two panes, a list of
\ctp{Class selector} lines above and the source below, with the same
middle-button menu on the list. You can read, edit and accept in it exactly as
in the full browser.

\section{The code pane menu}

Every text pane in the system shares one editor menu
(Section~\ref{sec:texteditor}). In a browser's code pane it carries three
more:

\begin{description}
\item[\ctp{format}] Pretty-print the method: the system's own indentation,
  computed by parsing it. It refuses --- and the window flashes --- if the
  text has unaccepted edits, because it works from the compiled method and not
  from what is on the screen. Hold the left shift key for the tighter
  \ctp{miniFormat}.
\item[\ctp{spawn}] A browser on just this method's text.
\item[\ctp{explain}] Select a token and choose this; an explanation of what
  that token \emph{is} --- temporary, instance variable, global, selector,
  class --- is inserted after it, as text, for you to read and then delete.
  The one piece of built-in documentation in the system.
\end{description}

One more editor item behaves differently here than its one-line description
suggests: \ct{again} repeats the last replacement at the next occurrence, and
with the left shift key held it makes the replacement throughout the whole
pane. That is the system's find-and-replace-all, and nothing announces it.

\section{Four questions and where to click}

\begin{description}
\item[``What does this message do?''] Select the selector in any code pane,
  or find it in a message list, and choose \ctp{implementors}.
\item[``Can I change this method?''] \ctp{senders}. Read them. There is no
  other answer.
\item[``What is this class for?''] Select it and choose \ctp{comment}, then
  \ctp{hierarchy} --- what it inherits from is usually a better description
  than its own methods.
\item[``Where did this variable come from?''] \ctp{inst var refs} on the
  class, and choose the variable. Every method that touches it, in one window.
\end{description}

And when the window is not the answer, Section~\ref{sec:withoutbrowser} has
the same questions asked from code and from the command line.
